\section{Test 1: Els nombres naturals}

\iniciTest{tA}{b,c,a,d,b,c,d,b,a,c,b,a}

\begin{enumerate}

\pregunta Quin axioma garanteix que \(0\) no és successor de cap nombre natural?
\begin{enumerate}
\opcio{a}{La injectivitat de \(S\).}
\opcio{b}{L'axioma que afirma que no existeix cap \(n\in\N\) tal que \(S(n)=0\).}
\opcio{c}{La commutativitat de la suma.}
\opcio{d}{La definició del producte.}
\end{enumerate}
\feedback

\pregunta La suma \(a+b\) es defineix recursivament per:
\begin{enumerate}
\opcio{a}{\(a+0=a\) i \(a+S(b)=S(a+b)\).}
\opcio{b}{\(a+0=0\) i \(a+S(b)=a+b\).}
\opcio{c}{\(a+b=b+a\), que és una definició.}
\opcio{d}{\(a+b=a\cdot b\).}
\end{enumerate}
\feedback

\pregunta El producte \(a\cdot b\) es defineix recursivament per:
\begin{enumerate}
\opcio{a}{\(a\cdot 0=a\) i \(a\cdot S(b)=a\cdot b\).}
\opcio{b}{\(a\cdot 0=1\) i \(a\cdot S(b)=S(a\cdot b)\).}
\opcio{c}{\(a\cdot b=b\cdot a\), que és una definició.}
\opcio{d}{\(a\cdot 0=0\) i \(a\cdot S(b)=a\cdot b+a\).}
\end{enumerate}
\feedback

\pregunta Segons la definició d'ordre en \(\N\), \(a\leq b\) vol dir:
\begin{enumerate}
\opcio{a}{Existeix \(n\in\N\) tal que \(a=b+n\).}
\opcio{b}{Existeix \(n\in\N\) tal que \(b=a+n\).}
\opcio{c}{\(a\) divideix \(b\).}
\opcio{d}{\(a\) i \(b\) són primers entre si.}
\end{enumerate}
\feedback

\pregunta La tricotomia de l'ordre afirma que, per a tots \(a,b\in\N\):
\begin{enumerate}
\opcio{a}{Sempre es compleixen alhora \(a<b\), \(a=b\) i \(b<a\).}
\opcio{b}{Sempre es compleix \(a<b\).}
\opcio{c}{Es compleix exactament una de les tres possibilitats: \(a<b\), \(a=b\) o \(b<a\).}
\opcio{d}{Cap nombre natural pot ser comparat amb un altre.}
\end{enumerate}
\feedback

\pregunta Si \(a\mid b\), amb \(a,b\in\N\) i \(a\neq 0\), això vol dir que:
\begin{enumerate}
\opcio{a}{\(a=b\).}
\opcio{b}{\(b\mid a\).}
\opcio{c}{\(a\) és primer.}
\opcio{d}{Existeix \(q\in\N\) tal que \(b=a\cdot q\).}
\end{enumerate}
\feedback

\pregunta Quin és \(\operatorname{mcd}(84,90)\)?
\begin{enumerate}
\opcio{a}{\(6\).}
\opcio{b}{\(12\).}
\opcio{c}{\(18\).}
\opcio{d}{\(1260\).}
\end{enumerate}
\feedback

\pregunta Quin és \(\operatorname{mcm}(4,6)\)?
\begin{enumerate}
\opcio{a}{\(4\).}
\opcio{b}{\(6\).}
\opcio{c}{\(12\).}
\opcio{d}{\(24\).}
\end{enumerate}
\feedback

\pregunta Quin dels nombres següents és primer?
\begin{enumerate}
\opcio{a}{\(21\).}
\opcio{b}{\(23\).}
\opcio{c}{\(25\).}
\opcio{d}{\(27\).}
\end{enumerate}
\feedback

\pregunta Quan és especialment útil el principi d'inducció forta?
\begin{enumerate}
\opcio{a}{Quan per demostrar el cas \(n\) convé usar tots els casos anteriors.}
\opcio{b}{Quan no es vol usar cap hipòtesi d'inducció.}
\opcio{c}{Quan es treballa només amb nombres enters negatius.}
\opcio{d}{Quan es vol evitar definir l'ordre en \(\N\).}
\end{enumerate}
\feedback

\end{enumerate}
